\section{Trace 怎样整理二次型微分}
\label{sec:10-Matrix-Calculus/03-Trace-Calculus/01}

训练脚本里加了一项自定义正则，惩罚特征之间的相关：\(\Omega(W)=\trace(W^{\trans}\Sigma W)\)，其中 \(\Sigma\) 是一个由数据估出来的常矩阵。框架不认识这个式子，反向要自己写。手边的记忆是``二次型的梯度是 \(2\Sigma W\)''，写进去，梯度检查报错，数值梯度和解析梯度差了一个有规律的量。

把规模缩到能心算的地步就看清楚了。取 \(A=\begin{bmatrix}0&2\\0&0\end{bmatrix}\)，\(f(x)=x^{\trans}Ax=2x_1x_2\)，真实梯度是 \((2x_2,\,2x_1)^{\trans}\)。按 \(2Ax\) 算得到 \((4x_2,\,0)^{\trans}\)，两者差 \((-2x_2,\,2x_1)^{\trans}\)。\(2Ax\) 这条公式要求 \(A\) 对称，而估出来的 \(\Sigma\) 若做过某种非对称的平滑或裁剪，就不满足。二阶矩矩阵通常对称，一旦经手过别的算子，这个前提就得重新确认。

\(2\times2\) 的错误可以靠展开元素抓出来。变量升级成矩阵之后就不行了——\(f(W)=\trace(W^{\trans}\Sigma W B)\) 展开成元素是四重求和，没人愿意核对。推导阶段必须一次做对，而做对的工具是 trace。

\subsection{trace 把乘积摆成一个可以旋转的环}

trace 有三条性质：对加法线性，对转置不变，以及在形状相容时允许因子整体循环移位，\(\trace(ABC)=\trace(BCA)=\trace(CAB)\)。第三条是主力。它说的是矩阵乘积在 trace 里没有固定的``开头''，四个因子首尾相接连成一个环，从环上任何一处切开读出来，值都一样。

这个环有一个特殊的位置，姑且叫做读数口。上一单元把梯度定义成使 \(df=\ip{G}{dX}=\trace(G^{\trans}dX)\) 成立的那个 \(G\)（见\hyperref[sec:10-Matrix-Calculus/02-Differentials/01]{``用微分与 Frobenius 内积读取梯度''}）。这个标准形要求 \(dX\) 站在紧挨读数口的位置，它左边剩下的那一串就是 \(G^{\trans}\)。于是矩阵求导的整理工作有了一个明确的目标：把环转到 \(dX\) 落在读数口。

\begin{center}
\begin{tikzpicture}[font=\small,>=stealth]
  \draw[black!30,thick] (0,0) circle (1.3);
  \draw[->,black!50,thick] (78:1.34) arc (78:56:1.34);
  \draw[dashed,black!60] (90:0.98) -- (90:1.68);
  \node[font=\footnotesize,black!60,anchor=south] at (90:1.72) {读数口};
  \node[draw,black!45,fill=white,rounded corners,inner sep=2.5pt] at (45:1.3) {\(X^{\trans}\)};
  \node[draw,black!45,fill=white,rounded corners,inner sep=2.5pt] at (-45:1.3) {\(A\)};
  \node[draw,blue!60!black,fill=white,rounded corners,inner sep=2.5pt] at (-135:1.3) {\(dX\)};
  \node[draw,black!45,fill=white,rounded corners,inner sep=2.5pt] at (135:1.3) {\(B\)};
  \node[font=\footnotesize,align=center] at (0,-2.35) {\(\trace(X^{\trans}A\,dX\,B)\)\\[1pt] \(dX\) 卡在中间};
  \begin{scope}[shift={(5.6,0)}]
    \draw[black!30,thick] (0,0) circle (1.3);
    \draw[->,black!50,thick] (78:1.34) arc (78:56:1.34);
    \draw[dashed,black!60] (90:0.98) -- (90:1.68);
    \node[font=\footnotesize,black!60,anchor=south] at (90:1.72) {读数口};
    \draw[red!65!black,very thick] (-8:1.62) arc (-8:-192:1.62);
    \node[draw,blue!60!black,fill=white,rounded corners,inner sep=2.5pt] at (45:1.3) {\(dX\)};
    \node[draw,black!45,fill=white,rounded corners,inner sep=2.5pt] at (-45:1.3) {\(B\)};
    \node[draw,black!45,fill=white,rounded corners,inner sep=2.5pt] at (-135:1.3) {\(X^{\trans}\)};
    \node[draw,black!45,fill=white,rounded corners,inner sep=2.5pt] at (135:1.3) {\(A\)};
    \node[font=\footnotesize,red!65!black,anchor=north] at (-90:1.72) {\(G^{\trans}=BX^{\trans}A\)};
    \node[font=\footnotesize,align=center] at (0,-2.35) {转一格之后\\[1pt] 梯度直接读出};
  \end{scope}
\end{tikzpicture}

\emph{同一个 trace 的两种读法：把环转到 \(dX\) 紧挨读数口，它左边剩下的一串就是梯度的转置。}
\end{center}

图里右边那一段红弧就是要读的东西。\(G^{\trans}=BX^{\trans}A\)，于是 \(G=A^{\trans}XB^{\trans}\)。整个推导没有触碰任何一个元素，只是把乘积在环上转了一格。

从最简单的情形建立读写节奏。\(f(X)=\trace(A^{\trans}X)\) 的微分是 \(df=\trace(A^{\trans}dX)\)，\(dX\) 已经在读数口，对照标准形直接得 \(\nabla_X f=A\)。再进一步，\(f(X)=\trace(AXB)\) 的微分是 \(\trace(A\,dX\,B)\)，转一格成 \(\trace(BA\,dX)\)，于是 \(G^{\trans}=BA\)，\(\nabla_X f=(BA)^{\trans}=A^{\trans}B^{\trans}\)。最后那个转置最容易漏。\(dX\) 左边的那一串是梯度的转置，不是梯度。

\subsection{向量二次型里反对称部分自己消失}

回到开头。\(f(x)=x^{\trans}Ax\) 的微分是 \(df=(dx)^{\trans}Ax+x^{\trans}A\,dx\)。第一项是标量，转置它不改变值：\((dx)^{\trans}Ax=x^{\trans}A^{\trans}dx\)。合并两项，

\[
df=x^{\trans}(A+A^{\trans})dx
=\bigl[(A+A^{\trans})x\bigr]^{\trans}dx,
\qquad
\nabla_x f=(A+A^{\trans})x.
\]

\(A\) 对称时它收缩成 \(2Ax\)。促成收缩的是对称性，不是``二次型''三个字，这两个条件值得分开记。

反对称部分为什么不出现在函数值里，也有一行答案。若 \(W^{\trans}=-W\)，则 \(x^{\trans}Wx=(x^{\trans}Wx)^{\trans}=x^{\trans}W^{\trans}x=-x^{\trans}Wx\)，故恒为零。所以 \(x^{\trans}Ax\) 只看得见 \(A\) 的对称部分 \(\tfrac12(A+A^{\trans})\)。梯度公式 \((A+A^{\trans})x\) 与这件事一致：它等于两倍的对称部分乘 \(x\)。

带一次项的完整形式 \(f(x)=\tfrac12x^{\trans}Ax+b^{\trans}x+c\) 的梯度是 \(\tfrac12(A+A^{\trans})x+b\)。当 \(A\) 是 Hessian 或协方差矩阵时它本来就对称，这一步是恒等变形；当 \(A\) 来路不明时，这一步救回一个梯度检查。

\subsection{变量出现几次，微分就有几条路径}

矩阵二次型 \(f(X)=\trace(X^{\trans}AXB)\) 里，\(X\) 出现两次。微分是线性运算，它同时击中两处，产生两项：

\[
df=\trace\bigl((dX)^{\trans}AXB\bigr)+\trace\bigl(X^{\trans}A(dX)B\bigr).
\]

第一项的 \(dX\) 带着转置。用 \(\trace(M)=\trace(M^{\trans})\) 把整项翻过来，\(\trace\bigl((dX)^{\trans}AXB\bigr)=\trace\bigl((AXB)^{\trans}dX\bigr)\)，读出贡献 \(AXB\)。第二项就是上面那张图：转一格得 \(\trace\bigl(BX^{\trans}A\,dX\bigr)\)，读出贡献 \(A^{\trans}XB^{\trans}\)。相加，

\[
\nabla_X f=AXB+A^{\trans}XB^{\trans}.
\]

\(A,B\) 都对称时它收缩成 \(2AXB\)。漏掉一条路径，得到的是这个结果的一半，而且梯度检查会显示一个稳定的比例偏差——一半的梯度是最难察觉的一种错误，因为方向常常仍然大致正确，只是步长系统性地偏小。

开头那个正则项现在可以结案了。\(\Omega(W)=\trace(W^{\trans}\Sigma W)\) 是上式取 \(B=I\)、\(A=\Sigma\) 的特例，梯度为 \((\Sigma+\Sigma^{\trans})W\)。\(\Sigma\) 对称时确为 \(2\Sigma W\)，否则不是。

\subsection{Frobenius 损失的梯度是残差被两端拉回}

Frobenius 范数平方是元素平方和，写成 trace 就是 \(\norm{X}_F^2=\trace(X^{\trans}X)\)。它的微分有两条路径，各贡献 \(\trace(X^{\trans}dX)\)，于是 \(d\norm{X}_F^2=2\trace(X^{\trans}dX)\)，\(\nabla_X\tfrac12\norm{X}_F^2=X\)。

把它套进一个带线性算子的损失：

\[
f(X)=\tfrac12\norm{AXB-C}_F^2.
\]

记残差 \(R=AXB-C\)。外层给出 \(df=\trace(R^{\trans}dR)\)，内层给出 \(dR=A\,dX\,B\)，两者复合：

\[
df=\trace(R^{\trans}A\,dX\,B)=\trace(BR^{\trans}A\,dX)
=\trace\bigl((A^{\trans}RB^{\trans})^{\trans}dX\bigr),
\]

读出 \(\nabla_X f=A^{\trans}(AXB-C)B^{\trans}\)。

这个式子在图像恢复里有具体身份。设 \(X\) 是待恢复的高分辨率图像，\(A\) 对列做模糊与下采样，\(B\) 对行做同样的事，\(C\) 是观测到的低分辨率图像，前向算子就是 \(X\mapsto AXB\)。梯度告诉你：先算残差，再用 \(A^{\trans}\) 和 \(B^{\trans}\) 把它送回高分辨率空间。

送回去的这一步值得说清楚，因为直觉容易把它读成``把模糊倒放''。\(A^{\trans}\) 并不是 \(A\) 的逆，下采样丢掉的方向落在算子的零空间里，那里的分量无论取什么值都产生同一个观测，残差对它们一无所知。伴随只能把误差投影回它够得着的子空间；零空间中的细节要靠正则化或先验补，不能指望反投影变出来。用抽象一点的话说，前向映射 \(X\mapsto AXB\) 是线性算子，它在 Frobenius 内积下的伴随是 \(G\mapsto A^{\trans}GB^{\trans}\)，梯度公式里残差两端各挨一个转置，就是在执行这个伴随。这正是\hyperref[sec:10-Matrix-Calculus/02-Differentials/02]{``链式法则就是线性映射复合''}中``反向传的是伴随''那句话的矩阵版本。

\subsection{动笔前的五项对账}

trace 是整理工具，不是推导的起点。推导的起点是微分，trace 负责把微分摆成能读的形状。下面几件事在动笔前确认一遍，能省掉大量返工。

先对形状。每个乘积的相邻维度必须匹配，一处不符，后面所有循环移位都失去意义；\(\trace(ABC)\) 要求 \(CA\) 也是方阵，这个隐含条件常被忽略。再点变量。哪些矩阵随参数变，哪些是常量，常量的微分为零，把它们先标出来能让式子短一半。然后数出现次数。变量在表达式中出现几次，微分就有几项，转置形式的出现（如 \(X^{\trans}\)）也算一次，且它带来的那一项需要先整体转置再读。接着看约束。若变量本身被限制为对称、对角、正交或低秩，无约束地读出的 \(G\) 还要投影到约束的切空间上，对称化 \(\tfrac12(G+G^{\trans})\) 是最常见的一种投影，直接拿未投影的梯度去更新会立刻破坏结构。最后守住循环的边界。\(\trace(ABCD)\) 只有四种合法读法，即 \(BCDA\)、\(CDAB\)、\(DABC\) 与它自己，把它写成 \(\trace(ACBD)\) 是把循环误当成了交换。

复矩阵不在本单元的默认设定里。那里要区分转置与共轭转置，梯度的定义也要换成 Wirtinger 导数，实矩阵的公式不能直接套。

到这里，被求导的对象都是 \(X\) 的多项式，每一步都是加法和乘法。下一节换一种出现方式：\(X\) 在表达式里以 \(X^{-1}\) 的身份登场。公式本身只有一行，麻烦的是它落到代码里的写法——那一行公式里的三个 \(X^{-1}\)，一个都不该显式算出来。
